\documentclass[UTF8]{ctexart}
\usepackage{amsmath,amssymb,bm}
\usepackage{enumitem}
\usepackage[hidelinks]{hyperref}
\usepackage[a4paper,margin=2.2cm]{geometry}
\setlist[itemize]{leftmargin=*, itemsep=2pt, topsep=2pt}
\setlist[enumerate]{leftmargin=*, itemsep=2pt, topsep=2pt}

\begin{document}

\section*{WS-OVVIS 技术路线大纲（v12 主线极简版）}

\subsection*{0. 名称与定位}
本文研究 \textbf{Weakly-supervised Open-Vocabulary Video Instance Segmentation}（WS-OVVIS）的一个严格设定：
\begin{itemize}
\item 训练阶段仅提供 \textbf{clip-level} 正证据标签集合（positive label set），且该集合系统性不完备；
\item 不使用像素级类别标注、实例级类别标注、box-level 实例语义监督，也不使用额外图文配对训练数据；
\item 允许使用 class-agnostic 的伪实例结构监督，但该结构监督本身不携带开放词汇语义；
\item 本文的核心主张不是生成质量更优的 tracklet，而是：\textbf{给定一个冻结的视频原生 class-agnostic basis，如何构建可复用的语义载体 bank，并在不完备 clip-level 正证据下完成 open-world 语义归因，且在测试时支持 bag-free 的 text-conditioned open-vocabulary inference。}
\end{itemize}

\par\noindent\hrulefill\par

\subsection*{1. 任务定义与边界说明}

\subsubsection*{1.1 输入、输出与监督边界}
训练输入：
\begin{itemize}
\item 一个原始视频 clip，记为 $v$；
\item 该 clip 的不完备正证据集合 $Y'(v)$，其中 $Y'(v)\subseteq Y^*(v)$，$Y^*(v)$ 为真实类别集合；
\item class-agnostic 的伪实例结构监督。本文固定沿用 \textbf{MaskCut + ImageCut2Video} 生成伪 video labels 的方式，不改变其标签生成逻辑。
\end{itemize}

测试输入：
\begin{itemize}
\item 视频 clip $v$；
\item 一个文本类别集合 $\mathcal C = \{c_1,\dots,c_K\}$，默认取当前 benchmark 的评测类别全集；
\item \textbf{测试时不输入} $Y'(v)$，也不依赖任何视频内正证据集合做候选约束。
\end{itemize}

输出：
\begin{itemize}
\item clip 内所有预测实例的 mask 序列；
\item clip 内每条 query trajectory 的开放词汇类别预测；
\item 可选保留 unknown foreground / reject 通道用于诊断，但标准 AP 仍按 top-1 已知类别口径汇报。
\end{itemize}

\subsubsection*{1.2 与经典 OVVIS 的关系}
\begin{itemize}
\item 经典 OVVIS 通常在实例级强监督下直接学习 ``实例$\leftrightarrow$类别'' 关系；
\item 本文仅在训练时获得不完备 clip-level 正证据，因此训练问题改写为 \textbf{bag-constrained open-world attribution}；
\item 为尽量对齐经典 OVVIS 的能力边界，测试阶段要求采用 \textbf{bag-free text-conditioned inference}，即不再依赖 $Y'(v)$ 做推理候选约束。
\end{itemize}

\subsubsection*{1.3 open-vocabulary 的主张边界}
\begin{itemize}
\item 主文的核心定义为：\textbf{benchmark-scoped open-vocabulary VIS}，即类别通过文本原型定义，测试时允许识别训练未见的 benchmark novel categories；
\item 主文正式采用 LV-VIS 官方 open-vocabulary 口径：训练期弱监督正证据仅暴露训练侧可见类别，测试期在完整 base+novel 词表上做 bag-free inference；
\item 方法在形式上支持 bag-free 的文本条件推理接口，因此比纯 bag-constrained 分类更接近经典 OVVIS；
\item 主文不把 unrestricted free-form text grounding 作为核心结论，也不把任意外部开放词表识别作为成立条件。
\end{itemize}

\par\noindent\hrulefill\par

\subsection*{2. 方法总览：主线极简闭环}
整体技术路线压缩为一条最小闭环：
\begin{itemize}
\item 保持 \textbf{MaskCut + ImageCut2Video} 不变，继续作为 class-agnostic pseudo video supervision 的生成方式；
\item 将 VideoCutLER 中原有的 \textbf{R50-VideoMask2Former} 替换为 \textbf{冻结 DINOv2 backbone + 轻量 adapter + VideoMask2Former decoder} 的 video-native learner；
\item 在 LV-VIS 域内沿用原有伪标签生成逻辑，训练新的 class-agnostic VideoMask2Former，使其直接输出视频原生的 query trajectories；
\item \textbf{主线中不再进行额外 stitching、tracklet refinement 或结构修复}，而是将每个 query trajectory 直接视为一个 atomic basis unit；
\item 对每个 atomic basis unit 提取 masked DINO pooling feature，构建可追溯、可复用的语义载体 bank；
\item 在语义载体 bank 上构建 seen 类视觉原型，并学习强约束的 class-level text-to-prototype map；
\item 训练时在 $Y'(v)+bg+unk$ 上进行 coverage-aware open-world attribution，测试时丢弃 bag 约束，直接进行 bag-free 的全词表文本条件推理。
\end{itemize}

\par\noindent\hrulefill\par

\subsection*{3. class-agnostic 伪视频监督：MaskCut + ImageCut2Video（保持不变）}

\subsubsection*{3.1 MaskCut：单图对象性伪 mask 生成}
\begin{itemize}
\item 对训练图像池中的每张图像，使用 MaskCut 生成 class-agnostic 的单图多对象伪 mask；
\item 该阶段只负责提供 ``对象性'' 结构先验，不提供语义类别标签；
\item 对每张图最多保留有限个对象 mask，并做边界细化与基础质量过滤，导出二值 mask、bbox 与基础 score。
\end{itemize}

\subsubsection*{3.2 ImageCut2Video：由单图伪 mask 合成伪视频轨迹}
\begin{itemize}
\item 继续沿用 ImageCut2Video，把单图伪对象转换为带时间维的伪 video labels；
\item 固定保留两类轨迹：\textbf{静止对象轨迹} 与 \textbf{移动对象轨迹}；
\item 移动对象通过随机缩放、平移、旋转、光照增强等简单几何与外观变换进行合成；
\item 输出监督形式保持为 class-agnostic 的伪 mask trajectories，不改变原方法关于标签生成的基本逻辑。
\end{itemize}

\subsubsection*{3.3 在本文中的角色}
\begin{itemize}
\item 这一阶段的唯一职责是提供 \textbf{class-agnostic pseudo video supervision}；
\item 它不负责开放词汇语义，也不直接为后续 prototype 或文本桥接提供语义信息；
\item 在本文主线中，MaskCut + ImageCut2Video 被视为固定的伪监督生成器，而不是需要继续重写的核心创新点。
\end{itemize}

\par\noindent\hrulefill\par

\subsection*{4. DINOv2-VideoMask2Former：新的 video-native basis generator}

\subsubsection*{4.1 设计目标}
新的 video-native 模块直接作为整条方法的 class-agnostic basis generator，整体结构改为：
\begin{center}
\textbf{pseudo video labels $\rightarrow$ DINOv2-VideoMask2Former $\rightarrow$ query trajectories}
\end{center}
其目标不是输出开放词汇类别，而是直接学习更强的 class-agnostic 视频实例结构。

\subsubsection*{4.2 Backbone 替换：冻结 DINOv2 + 轻量 adapter}
\begin{itemize}
\item 将原 VideoCutLER 中 VideoMask2Former 的 \textbf{R50 backbone} 替换为 \textbf{冻结 DINOv2 backbone}；
\item 为适配 VideoMask2Former 所需的多尺度输入，引入轻量 \textbf{ViT-style adapter / feature adapter}，把 DINOv2 的单尺度特征转化为供 pixel decoder 使用的多尺度特征；
\item 冻结 DINOv2 主干参数，仅训练 adapter、pixel decoder、transformer decoder 以及 mask/class heads；
\item 该设计的出发点是：利用 DINOv2 更强的外观表征与对象性先验，同时避免在 noisy pseudo supervision 下对大模型主干做不稳定的端到端更新。
\end{itemize}

\subsubsection*{4.3 VideoMask2Former 保留的部分}
\begin{itemize}
\item 保留 VideoMask2Former 的视频 mask decoder、query-based set prediction 和 3D mask trajectory 预测方式；
\item 继续采用 class-agnostic 的 ``前景对象 vs no-object'' 输出设置，不在该模块中引入开放词汇类别监督；
\item 继续使用视频级 mask trajectory 作为模型原生输出，而不是在该阶段显式学习文本对齐或类别映射。
\end{itemize}

\subsubsection*{4.4 训练域与训练目标}
\begin{itemize}
\item 在 LV-VIS 训练域内，沿用 3.1--3.2 的伪视频监督生成方式，训练新的 DINOv2-VideoMask2Former；
\item 训练目标仍然是 class-agnostic 的视频实例学习：学习 mask trajectory、跨帧 query 一致性以及前景对象的时序结构；
\item 这一阶段的输出被视为冻结的视频原生 basis，不在主线中再追加额外的结构 refinement 模块。
\end{itemize}

\par\noindent\hrulefill\par

\subsection*{5. Atomic basis：直接使用 query trajectories}

\subsubsection*{5.1 主线定义}
在本文主线中，每个 VideoMask2Former query trajectory 直接定义为一个 \textbf{atomic basis unit}，记为 $q$。主线中不再做额外的 cross-query stitching、tracklet refinement 或结构修复。

\subsubsection*{5.2 Atomic basis 的对象接口}
每个 atomic basis unit 统一表示为：
\[
 q = (\mathrm{query\_id},\ \{M_t\},\ \{b_t\},\ \mathcal T_q,\ s_q,\ \mathcal S_q),
\]
其中分别表示：
\begin{itemize}
\item query ID；
\item 跨帧 mask 序列 $\{M_t\}$；
\item bbox 序列 $\{b_t\}$；
\item 时间范围 $\mathcal T_q$；
\item query-level 置信度 $s_q$；
\item 来源模型与版本信息 $\mathcal S_q$，用于后续追溯。
\end{itemize}

\subsubsection*{5.3 选择这一主线的原因}
\begin{itemize}
\item 本文的核心主张是语义归因，而不是结构优化；
\item 因此前端 video-native 模块只需提供一个冻结的、class-agnostic 的 video-native basis；
\item query trajectory 已经是 VideoMask2Former 的原生输出对象，直接使用可使主线更干净、更稳定，也避免把论文贡献重新拉回到 tracklet 生成质量上。
\end{itemize}

\par\noindent\hrulefill\par

\subsection*{6. 语义载体 bank：masked DINO pooling 作为主要语义载体}

\subsubsection*{6.1 主语义载体定义}
对每个 atomic basis unit $q$，设其可见帧集合为 $\mathcal V(q)$。对每个帧 $t\in \mathcal V(q)$，用 query mask 在 DINOv2 特征图上做 masked pooling，得到帧级语义特征：
\[
 f_{q,t} = \mathrm{Pool}(F_t^{\mathrm{DINO}}, M_{q,t}).
\]
随后在可见帧上按可见性与基础质量加权聚合，得到 trajectory-level 语义载体：
\[
 z_q = \frac{\sum_{t\in \mathcal V(q)} w_{q,t}^{d} f_{q,t}}{\sum_{t\in \mathcal V(q)} w_{q,t}^{d}}.
\]
本文将 $z_q$ 视为后续开放词汇语义建模的 \textbf{唯一主语义载体}。

\subsubsection*{6.2 为什么不用 query logits 作为主语义载体}
\begin{itemize}
\item VideoMask2Former 的 class head 只承担前景对象 vs no-object 判别，不承担开放词汇语义；
\item query logits 更像一次性分类信号，而 $z_q$ 是可持久化、可复用的视觉表征；
\item 因此，后续原型构建、文本映射和开放词汇归因全部建立在 $z_q$ 上，而不是建立在 VideoMask2Former 的 class logits 上。
\end{itemize}

\subsubsection*{6.3 可追溯、可复用的语义载体 bank}
本文显式维护一个持久化的语义载体 bank，而不是仅做一次性在线判断。主线中至少保存以下两层：

\paragraph{(1) Query-Trajectory Bank}
每个条目对应一个 atomic basis unit：
\begin{itemize}
\item query ID、video ID、帧索引、mask 序列、bbox 序列；
\item query score；
\item 来源模型版本与生成配置。
\end{itemize}
该层保证每个后续语义判断都能反查到其结构来源。

\paragraph{(2) Semantic Carrier Bank}
每个条目对应一个 trajectory-level 语义载体：
\begin{itemize}
\item query ID；
\item trajectory-level DINO embedding $z_q$；
\item pooling 权重与聚合配置；
\item 当前语义状态（未归因、已归因、unknown foreground 等）。
\end{itemize}
一旦 $z_q$ 被物化保存，后续 prototype 更新、文本映射修正和语义重赋值都可在 \textbf{不重新运行 VideoMask2Former} 的前提下完成。

\subsubsection*{6.4 在整条主线中的作用}
\begin{itemize}
\item Query-Trajectory Bank 提供可追溯的结构对象索引；
\item Semantic Carrier Bank 提供后续所有语义模块的直接输入；
\item 因此，本文后续不再把 ``track feature'' 定义为经过额外结构转换后的对象特征，而直接把每个 query trajectory 的 $z_q$ 视为语义层的基本单元。
\end{itemize}

\par\noindent\hrulefill\par

\subsection*{7. 训练集内部类视觉原型与 class-level text map}

\subsubsection*{7.1 文本侧与视觉侧角色分离}
\begin{itemize}
\item 文本侧固定使用 \textbf{冻结 CLIP Text Encoder + 类名 only}；主文不默认使用 prompt ensemble、同义词扩展或 LLM 描述增强；
\item 视觉侧统一基于 6.1 的 trajectory-level DINO 语义载体 $z_q$ 构建视觉原型；
\item 文本侧不直接去回归视频实例，而是通过 class-level map 映射到视觉原型空间。
\end{itemize}

\subsubsection*{7.2 seen 类视觉原型}
\begin{itemize}
\item 在训练集内部，对每个 seen/base 类 $y$，在所有满足 $y\in Y'(v)$ 的 clip 中，依据 clip-level 正证据覆盖启发式选择 latent positive query trajectories；
\item 用这些 latent positive trajectories 的语义载体 $z_q$ 构建该类的视觉原型 $p_y$；
\item prototype bank 采用 ``离线初始化 + 训练中低频刷新'' 的策略，不把复杂 EM 轮数与频繁自训练写入主线。
\end{itemize}

\subsubsection*{7.3 class-level text-to-prototype map}
在此基础上，仅学习一个 \textbf{强约束的 class-level map} $A$，使文本嵌入 $t_y$ 能映射到视觉原型附近：
\[
\hat p_y = \mathrm{LN}(A\,\mathrm{LN}(t_y)).
\]
主实现采用 \textbf{单层线性映射 $A$ + 强 L2 正则}，并使用正交 Procrustes 作为初始化；不使用大 MLP，也不做 instance-level 自由文本桥接。

\par\noindent\hrulefill\par

\subsection*{8. 训练期：coverage-aware open-world attribution}

\subsubsection*{8.1 输入对象}
训练阶段的归因对象为：
\begin{itemize}
\item 当前 clip 的 semantic carrier set $Z(v)=\{z_q\}$；
\item 当前 clip 的 atomic basis set $Q(v)=\{q\}$；
\item 当前 clip 的正证据集合 $Y'(v)$；
\item mapped text prototypes $\hat P=\{\hat p_j\}$。
\end{itemize}

\subsubsection*{8.2 主线原则}
\begin{itemize}
\item 训练只在 \textbf{$Y'(v)+bg+unk$} 上进行归因：$Y'(v)$ 用于承接已观测到的正证据，$bg$ 用于吸收低质量或非前景实例，$unk$ 用于承接当前未被正证据解释、但又不应塌成背景的前景质量；
\item 核心目标不是让每个 query trajectory 必然落到某个已知类，而是让 ``已知正证据得到覆盖'' 与 ``hidden positives 不被压成背景'' 同时成立；
\item 因此主方法必须显式保留 \textbf{unknown foreground slack} 与 \textbf{coverage constraint}。
\end{itemize}

\subsubsection*{8.3 相似度与代价}
对 $j\in Y'(v)$，定义语义相似度与 logit：
\[
 s_{qj}=\cos(z_q,\hat p_j),\qquad l_{qj}=s_{qj}/T_{temp}.
\]
对语义列 $j\in Y'(v)$，定义代价 $C_{qj}=-l_{qj}$。

背景列与 unknown 列的代价仅由前景质量控制：
\[
 C_{q,bg}=c_{bg}-\lambda_{bg}(1-s_q),
\]
\[
 C_{q,unk}=c_{unk}-\lambda_{unk}s_q.
\]
其含义是：低质量的 query trajectory 更容易被解释为 bg，而高质量但无法由当前正证据解释的 query trajectory 允许流向 unk。

\subsubsection*{8.4 归因主方法}
主实现采用 coverage-aware entropic soft assignment，在 $Y'(v)\cup\{bg,unk\}$ 上求软分配矩阵 $P^{(v)}$：
\[
 \min_{P^{(v)}\ge 0}\ \langle C,P^{(v)}\rangle - \varepsilon H(P^{(v)})
 + \lambda_{\mathrm{row}}\,\mathrm{KL}(P^{(v)}\mathbf 1\|a^{(v)})
 + \lambda_{\mathrm{cov}}\sum_{y\in Y'(v)} \Big[m_y^{(v)}-\sum_q P^{(v)}_{qy}\Big]_+^2.
\]
其中 $a_q^{(v)}\propto s_q$ 并归一化，表示每个 query trajectory 需被解释的质量；coverage 下界采用统一低下界
\[
 m_y^{(v)}=\rho_{cov}\cdot \frac{1}{N_v}\sum_q a_q^{(v)},\qquad y\in Y'(v),
\]
仅用于防止正证据被完全漏掉，而不编码每类实例数先验。

\subsubsection*{8.5 训练注入}
\begin{itemize}
\item 仅当某 query trajectory 在非 bg/unk 列上的总质量超过阈值时，才把 $P^{(v)}$ 在非 bg/unk 列上按行归一化为软标签，并用它监督对应的语义分布；
\item 结构模块与语义模块在职责上解耦：前者负责 class-agnostic video-native basis，后者负责 bag-constrained open-world attribution；
\item 主线默认不做大范围 end-to-end joint fine-tuning，而是把优化集中在 class-level map、prototype 更新与 attribution 相关标度参数上。
\end{itemize}

\par\noindent\hrulefill\par

\subsection*{9. 测试期：bag-free text-conditioned open-vocabulary inference}
\begin{itemize}
\item 测试阶段不再使用 $Y'(v)$，也不再围绕视频内正证据构造候选集；
\item 给定视频与文本类别集合 $\mathcal C$，对每个 query trajectory 与全部文本类别直接计算
\[
 s_{q,c}=\cos(z_q,\hat p_c),\qquad c\in\mathcal C.
\]
\item 推理沿用与训练一致的 class-agnostic structure pipeline：先得到 query trajectories，再对每个 trajectory 计算语义载体 $z_q$，最后做全词表文本匹配；
\item 主结果默认输出 top-1 known class；unknown foreground 分支主要用于诊断高质量但低文本可解释性的 trajectories，不作为标准 AP 的强依赖后处理。
\end{itemize}

\par\noindent\hrulefill\par

\subsection*{10. 扩展与技术备选方案（非主线）}
以下内容不属于本文主线，仅作为当主线出现问题时的技术备选：

\subsubsection*{10.1 结构侧扩展}
\begin{itemize}
\item 对 query trajectories 做 conservative stitching；
\item 对局部缺帧做 temporal repair；
\item 对高度重叠 query 做 duplicate suppression；
\item 将 query trajectories 转换为更长的 track-like objects 再进入语义模块。
\end{itemize}

\subsubsection*{10.2 语义侧防御机制}
\begin{itemize}
\item quality-aware semantic weighting：按时长、面积、分数稳定性对 $z_q$ 加权；
\item duplicate-aware prototype update：同一视频内高重叠 query 不同时高权更新 prototype；
\item uncertain foreground filtering：对极短、极小、极低分 trajectories 延迟归入 prototype bank。
\end{itemize}

\subsubsection*{10.3 更强的可复用 bank}
\begin{itemize}
\item 保存 frame-level masked DINO features，而不仅保存 trajectory-level $z_q$；
\item 保存 keyframe-only pooled features 或 multi-view pooled variants；
\item 在不重新运行 VideoMask2Former 的前提下，支持离线重算 pooling、prototype 与 attribution。
\end{itemize}

\subsubsection*{10.4 更强的视频原生模块备选}
\begin{itemize}
\item 将 DINOv2-VideoMask2Former 替换为其他更强的视频原生 query-based module；
\item 或在前端后接专门的 tracker/refiner；
\item 这些方向属于进一步提升结构质量的备选分支，不属于本文主线主张。
\end{itemize}

\par\noindent\hrulefill\par

\subsection*{11. 本版技术路线的核心变化总结}
相较上一版技术路线，本版做了以下实质性改写：
\begin{itemize}
\item 完全删除独立的中间 refinement 承接器，不再以生成更优 tracklet 作为论文主张；
\item 保留 MaskCut + ImageCut2Video，不修改伪视频监督生成逻辑；
\item 将 VideoCutLER 中的 VideoMask2Former 主干由 R50 改为 \textbf{冻结 DINOv2 + 轻量 adapter}，作为新的 video-native basis generator；
\item 主线中直接把每个 VideoMask2Former query trajectory 视为 atomic basis unit；
\item 以后续 masked DINO pooling 语义载体 bank 作为主方法核心，而不是以后处理结构优化为核心；
\item 所有质量防御、duplicate 处理、temporal repair 与更强可复用 bank 均后移到扩展技术选项，不写入主线；
\item 整条主线现在变为：
\begin{center}
\textbf{MaskCut + ImageCut2Video $\rightarrow$ DINOv2-VideoMask2Former $\rightarrow$ Query Trajectories $\rightarrow$ Semantic Carrier Bank $\rightarrow$ Prototype + Text Map $\rightarrow$ Open-World Attribution $\rightarrow$ Bag-free Inference}
\end{center}
\end{itemize}


\par\noindent\hrulefill\par

\subsection*{12. 门控设计：工程闭环与科学验证}
本节在当前主线极简版技术路线之上，单独给出工程门控与科学门控。门控的设计原则如下：
\begin{itemize}
\item 门控服务于当前主线主张：\textbf{本文的核心不再是生成质量更优的 tracklet，而是给定冻结的视频原生 class-agnostic basis，完成可复用语义载体 bank 的构建与弱监督 open-world 语义归因。}
\item 因此，工程门控只检查主线是否真实闭环、核心工件是否可落地、是否可追溯与可复用；科学门控只检查当前论文主张是否成立，不再把结构优化本身作为前置成立条件。
\item 如无必要，主线参数保持原始 VideoCutLER 默认设置；涉及 teacher surface、confidence threshold、object mask threshold 与 overlap threshold 的修改，均不得依据已弃用的 SeqFormer 路线结果直接进入主线，而应在后续扩展分支中单独验证。
\item 对于 VideoCutLER 输出，主线默认保留原始阈值体系（例如官方配置中的 object-mask / overlap 相关阈值），不把 low-confidence 调参写入主线定义；只有当新的 DINOv2-VideoMask2Former 主线在原始阈值下未通过非劣门时，才允许开启阈值扫描扩展分支。
\end{itemize}

\subsubsection*{12.1 工程门控}

\paragraph{E0. 主线配置冻结门}
目的：冻结当前论文主线的结构定义，防止在工程推进中重新引入旧版 refinement 依赖或无根据的阈值调参。
\begin{itemize}
\item 固定 MaskCut + ImageCut2Video，不修改其伪视频监督生成逻辑；
\item 固定 DINOv2 + 轻量 adapter + VideoMask2Former 为唯一 video-native learner；
\item 固定主线直接使用 query trajectories 作为 atomic basis，不启用 stitching、refinement、duplicate suppression、cardinality control；
\item 固定保留 Query-Trajectory Bank 与 Semantic Carrier Bank；
\item 默认沿用 VideoCutLER 原始阈值体系，不将任何基于 SeqFormer 的 low-confidence teacher surface 写入主线。
\end{itemize}
通过条件：训练、导出、推理与评估脚本仅引用一份主线配置；任何阈值修改或结构增强均必须显式归入“扩展与技术备选方案”。

\paragraph{E1. DINOv2-VideoMask2Former 集成闭环门}
目的：确认新的视频原生 basis generator 已被正确接入并可以稳定训练与导出。
\begin{itemize}
\item 冻结 DINOv2 backbone，adapter、pixel decoder、transformer decoder、queries 与 mask/class heads 正常训练；
\item 训练过程中无 NaN、无梯度爆炸、无全局空预测塌缩；
\item 推理阶段能够稳定导出非空 query trajectory 集合。
\end{itemize}
通过条件：最小训练闭环与验证闭环均可完成，且输出包含 \texttt{query\_id / video\_id / masks / boxes / frame\_ids / score / model\_version} 等核心字段。

\paragraph{E2. Query-Trajectory Bank 物化门}
目的：确认结构来源层真实存在并可持久化，而不是临时内存对象。
\begin{itemize}
\item 每条 query trajectory 均有稳定主键；
\item 可从 bank 反查到视频、帧、mask、bbox、score 及模型版本；
\item 训练与测试两侧 schema 保持一致；
\item 支持增量写入与随机读取。
\end{itemize}
通过条件：随机抽样若干 query trajectory，可从任一后续语义结果反查回其结构来源。

\paragraph{E3. Semantic Carrier Bank 物化门}
目的：确认 $z_q$ 不是在线即弃的中间变量，而是可复用的持久化语义工件。
\begin{itemize}
\item 对每个 atomic basis unit 均能生成 trajectory-level 语义载体 $z_q$；
\item 保存 pooling 配置、query ID、语义状态与必要元数据；
\item 在不重新运行 VideoMask2Former 的前提下，可重新构造 prototype 与 attribution 输入。
\end{itemize}
通过条件：使用已存的 Semantic Carrier Bank，可独立重建 prototype bank 并重新执行 attribution，而无需重新推理前端视频模块。

\paragraph{E4. 主线端到端闭环门}
目的：确认当前极简主线真实成立。
主线闭环如下：
\begin{center}
\textbf{MaskCut + ImageCut2Video $\rightarrow$ DINOv2-VideoMask2Former $\rightarrow$ Query Trajectories $\rightarrow$ Query-Trajectory Bank $\rightarrow$ Semantic Carrier Bank $\rightarrow$ Prototype + Text Map $\rightarrow$ Open-World Attribution $\rightarrow$ Bag-free Inference}
\end{center}
通过条件：训练期 attribution loss 可正常计算；测试期在不给定 $Y'(v)$ 的前提下可完成 bag-free inference；且整个闭环不依赖任何扩展侧结构增强模块。

\paragraph{E5. 可复用性回放门}
目的：确认 bank 机制不是简单缓存，而是真正支持后续复用与重赋值的中间知识层。
通过条件：在固定 Query-Trajectory Bank 与 Semantic Carrier Bank 的前提下，不重新运行前端视频模块，也能够：
\begin{itemize}
\item 重新构造或刷新 seen prototype bank；
\item 重新执行 text map 与 attribution；
\item 重新输出每个 query trajectory 的语义标签与匹配结果。
\end{itemize}

\subsubsection*{12.2 科学门控}

\paragraph{S1. 新 basis generator 非劣门}
科学问题：将 VideoCutLER 中原有的 R50 backbone 替换为冻结 DINOv2 + adapter 后，新的 DINOv2-VideoMask2Former 作为 class-agnostic basis generator，是否至少不劣于原始 R50-VideoMask2Former 输出。
\begin{itemize}
\item 比较对象：\texttt{VC-R50 raw output @ original thresholds} 与 \texttt{VC-DINOv2 raw output @ original thresholds}；
\item 核心指标：\texttt{mean\_best\_iou}、\texttt{recall@0.5}、\texttt{fragmentation\_per\_gt\_instance}、\texttt{predicted\_tracks\_total}；
\item 通过原则：不要求全面显著优于 R50，但要求不出现明显性能退化与 coverage-side collapse。
\end{itemize}
若 S1 在原始阈值下失败，才允许进入阈值扫描扩展分支；主线默认不以调低阈值作为前提条件。

\paragraph{S2. DINO pooled feature 作为主语义载体成立门}
科学问题：后续语义建模是否应建立在 trajectory-level masked DINO pooled feature $z_q$ 上，而不是建立在 VideoMask2Former 的 query logits 上。
\begin{itemize}
\item 比较对象：query logits 路线 vs trajectory-level DINO pooled $z_q$ 路线；
\item 观察重点：seen prototype 的类内聚合稳定性、text-to-prototype map 的可学习性、attribution 的稳定性与测试期 bag-free inference 的可用性；
\item 通过原则：只要 $z_q$ 路线在下游语义指标或稳定性上明显更优，即认为“DINO pooled feature 是唯一主语义载体”的设计成立。
\end{itemize}

\paragraph{S3. 弱监督 open-world attribution 成立门}
科学问题：在只给定不完备 clip-level 正证据集合 $Y'(v)$ 的条件下，coverage-aware attribution 是否真正学会了“覆盖已观测正证据，同时不把 hidden positives 系统性压成背景”。
\begin{itemize}
\item 需要检查正证据列的覆盖率是否足够；
\item 需要检查高质量 query trajectories 是否更倾向流向 unknown foreground，而非被塌缩为背景；
\item 需要检查与无 unk / 无 coverage constraint 版本相比，主方法是否表现出更合理的前景质量分流。
\end{itemize}
通过原则：只要已观测正证据能够稳定获得 attribution mass，且 hidden positives 不被系统性压成背景，则 S3 成立。

\paragraph{S4. Bag-free inference 成立门}
科学问题：训练阶段使用 bag 约束后，测试时丢弃 bag 约束，系统是否仍能在全词表上进行可用的 text-conditioned open-vocabulary inference。
\begin{itemize}
\item 比较对象：bag-constrained inference（诊断上界）与 bag-free inference（正式主张）；
\item 观察重点：bag-free 设置下对 base+novel 词表的区分能力，以及相较 bag-constrained 版本是否出现不可接受的崩溃。
\end{itemize}
通过原则：允许 bag-free inference 相对 bag-constrained 有可接受退化，但不允许在去掉 bag 后完全失效。

\paragraph{S5. 可复用语义 bank 科学成立门}
科学问题：当前 Semantic Carrier Bank 是否真正具有后续利用价值，而非一次性中间缓存。
\begin{itemize}
\item 使用同一批已物化的 query trajectories 与 $z_q$，在不重新运行前端视频模块的条件下，重新构建 prototype、更新 text map、重新执行 attribution；
\item 检查 bank 驱动的重赋值是否与在线版本保持一致，且每个语义结果均可回查到 Query-Trajectory Bank 中的结构来源。
\end{itemize}
通过原则：若语义重赋值、prototype 重建与错误分析均可仅依赖已存 bank 完成，则 S5 成立。

\subsubsection*{12.3 失败后的扩展触发条件（非主线）}
以下触发条件仅用于指导后续扩展，不改变主线定义：
\begin{itemize}
\item 若 S1 失败：优先排查 adapter 设计、pixel decoder / transformer decoder 训练与 backbone 冻结策略；只有在原始阈值下持续不通过非劣门时，才允许开启阈值扫描扩展分支。
\item 若 S2 失败：优先排查 masked pooling 设计、帧权重、关键帧池化与 frame-level feature 物化策略，而不是立即引入结构 refinement。
\item 若 S3 失败：再考虑 quality-aware weighting、duplicate-aware prototype update、uncertain foreground filtering 等语义侧防御机制。
\item 若 S4 失败：优先排查 prototype 污染、text map 过弱或 train/test 接口不一致，而不是回退到 bag-constrained 推理。
\item 若 S5 失败：再考虑扩展为更强可复用 bank，例如显式保存 frame-level masked DINO features 或 keyframe-only pooled variants。
\end{itemize}

\end{document}
